%! Sinusoidal Functions — Unit 6, Lesson 6.6 — Guided Notes (Answer Key)
\documentclass[10pt]{article}
\usepackage{precalculus-article}
\usepackage{precalculus-key}
\newcommand{\TallMath}[1]{$\displaystyle #1\rule[-1.4em]{0pt}{3.2em}$}

\begin{document}

\pageheader{Unit 6, Lesson 6.6}{Guided Notes --- Answer Key}
\namedateperiod

\vspace{0.1in}

%-----------------------------------------------------------------------
% OBJECTIVE
%-----------------------------------------------------------------------
\begin{objectivebox}
\textbf{I will be able to\ldots}
\begin{itemize}
  \item Define a sinusoidal function and explain how cosine is related to sine.
  \item Determine the amplitude, midline, period, and frequency of a sinusoidal
        function from a verbal description or graph.
  \item Describe the symmetry of $y=\sin\theta$ and $y=\cos\theta$.
\end{itemize}
\end{objectivebox}

\vspace{0.08in}

%-----------------------------------------------------------------------
% VOCABULARY
%-----------------------------------------------------------------------
\begin{vocabbox}
\begin{itemize}
  \item \termblanklong{sinusoidal function}{\ans{any function formed by additive
        and multiplicative transformations of $f(\theta)=\sin\theta$}}
  \item \termblanklong{amplitude}{\ans{half the difference between the maximum
        and minimum values}}
  \item \termblanklong{midline}{\ans{the horizontal line $y =
        \tfrac{\text{max}+\text{min}}{2}$, halfway between max and min}}
  \item \termblanklong{period}{\ans{the smallest positive $T$ such that
        $f(\theta + T) = f(\theta)$ for all $\theta$}}
  \item \termblanklong{frequency}{\ans{$1/T$; the reciprocal of the period}}
  \item \termblanklong{odd function}{\ans{$f(-x) = -f(x)$; rotational symmetry
        about the origin}}
  \item \termblanklong{even function}{\ans{$f(-x) = f(x)$; reflective symmetry
        over the $y$-axis}}
\end{itemize}
\end{vocabbox}

\vspace{0.08in}

%-----------------------------------------------------------------------
% HOOK
%-----------------------------------------------------------------------
\begin{hookbox}
\textit{``A recording engineer notices that a sound wave has amplitude 0.3 and
frequency 440 Hz.''}

What does amplitude tell us about the wave?

\medskip
\ans{The wave oscillates 0.3 units above and below its center (midline),
controlling perceived loudness.}

\medskip
What does frequency tell us?  What is the period of one cycle?

\medskip
\ans{Frequency 440 Hz means 440 complete cycles per second.
Period $= 1/440 \approx 0.00227$ seconds.}
\end{hookbox}

\vspace{0.08in}

%-----------------------------------------------------------------------
% CHARACTERISTICS OF SINUSOIDAL FUNCTIONS
%-----------------------------------------------------------------------
\begin{notesbox}{Characteristics of Sinusoidal Functions --- Key}

\textbf{Amplitude}
\[
  \text{Amplitude} = \frac{\ans{max} - \ans{min}}{2}
  = \text{half the } \ans{vertical range}
\]

\textbf{Midline}
\[
  \text{Midline: } y = \frac{\ans{max} + \ans{min}}{2}
  = \text{average of } \ans{max and min}
\]

\textbf{Period} \quad The smallest positive \ans{value $T$} such that
$f(\theta + k) = f(\theta)$ for all $\theta$.

\textbf{Frequency}
\[
  \text{Frequency} = \frac{\ans{1}}{\text{period}}
  = \frac{1}{\ans{T}}
\]

\medskip
\textbf{For $f(\theta) = \sin\theta$:}

\renewcommand{\arraystretch}{1.8}
\begin{tabular}{@{}llll@{}}
  Amplitude $=$ \ans{1} &
  Midline $=$ \ans{$y=0$} &
  Period $=$ \ans{$2\pi$} &
  Frequency $=$ \ans{$1/(2\pi)$}
\end{tabular}

\end{notesbox}

\vspace{0.08in}

%-----------------------------------------------------------------------
% RELATIONSHIP BETWEEN SINE AND COSINE
%-----------------------------------------------------------------------
\begin{notesbox}{Relationship Between Sine and Cosine --- Key}

Complete the identity:
\[
  \cos\theta = \sin\!\left(\theta + \ans{$\pi/2$}\right)
\]

This means cosine is a \ans{horizontal shift (phase shift)} of the sine function.

Therefore cosine is also a \ans{sinusoidal} function.

\end{notesbox}

\vspace{0.08in}

%-----------------------------------------------------------------------
% SYMMETRY
%-----------------------------------------------------------------------
\begin{notesbox}{Symmetry of Sine and Cosine --- Key}

\begin{multicols}{2}
\textbf{Sine --- Odd Function}

$\sin(-\theta) = $ \ans{$-\sin\theta$}

Graphical meaning: the graph has \ans{rotational}
symmetry about the \ans{origin}.

\columnbreak

\textbf{Cosine --- Even Function}

$\cos(-\theta) = $ \ans{$\cos\theta$}

Graphical meaning: the graph has \ans{reflective}
symmetry over the \ans{$y$-axis}.
\end{multicols}

\end{notesbox}

\vspace{0.08in}

%-----------------------------------------------------------------------
% PRACTICE
%-----------------------------------------------------------------------
\begin{practicebox}
\textbf{Practice:} A sinusoidal function has maximum value $5$, minimum value $-1$,
and period $4\pi$.  Find the amplitude, midline, and frequency.

\medskip
Amplitude $= \dfrac{\ans{5} - (\ans{$-1$})}{2} = $ \ans{3}

\medskip
Midline: $y = \dfrac{\ans{5} + (\ans{$-1$})}{2} = $ \ans{$y = 2$}

\medskip
Frequency $= \dfrac{1}{\ans{$4\pi$}} = $ \ans{$\dfrac{1}{4\pi}$}
\end{practicebox}

\end{document}
